\documentclass[10pt]{article}

% Specify the margins and text size.
\setlength{\textwidth}{6.4in}
\setlength{\textheight}{9.5in}
\setlength{\oddsidemargin}{0pt}
\setlength{\evensidemargin}{0pt}
\setlength{\topmargin}{0pt}
\setlength{\hoffset}{.05in}
\setlength{\voffset}{-1in}

\setlength{\parskip}{5pt}
\setlength{\parindent}{0pt}

% Load some fonts and symbol packages
\usepackage{latexsym}
\usepackage{pifont}       % contains 'star' symbol for counterinsurgency handbook title
\usepackage{yfonts} 
\usepackage{amsmath}
\usepackage{amsfonts}

\usepackage{graphicx}     % actually, this is loaded by pstricks
\usepackage[T1]{fontenc}
\usepackage{ifthen}
\usepackage{pstricks,pst-grad,pst-text,pst-node,multido,pst-plot,calc,pst-3dplot}
%\usepackage[all]{xy}
%\usepackage{animate}

% The hyperref package inserts links into the pdf.
\definecolor{MyLinkColor}{rgb}{.1,.2,1}
\definecolor{MyCiteColor}{rgb}{.1,1,.2}
\definecolor{MyURLColor}{rgb}{.4,.4,.4}
\usepackage[backref=true,pagebackref=false,hyperindex,colorlinks=true,
  linkcolor=MyLinkColor,urlcolor=MyURLColor]{hyperref}


% The tweaklist package is something I found on the web.  It provides a simple interface
% for making changes to spacing used in the itemize and enumerate environments.  Comment
% this out if you don't care to use tweaklist.
\usepackage{tweaklist}
\renewcommand{\itemhook}{\setlength{\parskip}{2pt}\setlength{\parsep}%
{1pt}\setlength{\topsep}{0pt}\setlength{\itemsep}{0pt}}

\newcommand{\U}{\underline{\hspace{5pt}}}

\usepackage{listings}
\newcounter{EX}\setcounter{EX}{1}
\newcommand{\EXERCISE}{\arabic{EX}.\stepcounter{EX} }

\begin{document}
\pagestyle{empty}
\lstset{language=R, showspaces=false, showstringspaces=false}

\href{http://www.su.edu}{\includegraphics[height=1.75cm]{sulogo.eps}}
\vspace{-1.69cm}

{{\ }\hfill\small
\begin{tabular}{cl}
& Math 207\\
& Introduction to Statistics\\
\end{tabular}
}
\setlength{\baselineskip}{1.05\baselineskip}
\bigskip

\begin{center}
\textbf{\large  Quiz 1}
\end{center}
\medskip
\newcommand{\Z}{\hphantom{0}}

1. A fair die is rolled 10 times.  The number of dots showing on the rolls
was:  2, 1, 4, 4, 2, 4, 1, 3, 1, 5. On the back side of this page, 
draw a histogram of this data.
\vspace{.5in}

2. A hypothetical population distribution is shown in the table below.

{\ }\hspace{10pt}\begin{tabular}{cc}
       & Percent of\\
Age    & Population\\\hline
0--5   & 20\vphantom{\Large Y}\\  % width=5, height=4
\Z6--12  & 18\\  % width=6, height=3
13--18 & 10\\  % width=5, height=2
19--34 & 30\\  % width=15, height=2
35--45 & 10\\  % width=10, height=1
46--56 & \Z5\\ % width=10, height=1/2
57--71 & \Z7   % width=14, height=1/2
\end{tabular}
\vspace{.5in}

\hspace{5pt} a) Draw a histogram of the population distribution.  Label
both the axes.

\hspace{5pt} b) About what percent  of the population is between 31 and 33 years old?
\vspace{.5in}

%\hspace{5pt} c) About what percent  of the population is between 44 and 46 years old?
%\vspace{.5in}

3. (Hypothetical.)  In a clinical trial, data collection usually starts 
at \lq\lq baseline,\rq\rq\ when the subjects are recruited into the trial 
but before they are assigned to treatment or control.  Data collection continues 
until the end of followup.  Two clinical trials on prevention of heart attacks 
report baseline data on smoking, shown below.  In one of these trials, 
the randomization did not work.  Which one was it and and how can you tell?
\vspace{-12pt}
\newcommand{\Zc}{\hphantom{0,}}
\begin{center}
\begin{tabular}{cc}
\begin{tabular}{l}
\\
\\
\rput(-12pt,-3pt){(i)}
Treatment\vphantom{\Large Y}\\
Control\\[10pt]
\rput(-12pt,-3pt){(ii)}
Treatment\\
Control
\end{tabular}
&
\begin{tabular}{cc}
Number of & Percent\\
persons & who smoked\\\hline
1,012  & 49.3\%\vphantom{\Large Y}\\
\Zc997  & 69.0\%\\[10pt]
\Zc995  & 59.3\%\\
1,017   & 59.0\%
\end{tabular}
\end{tabular}
\end{center}
\vfill
\eject


\href{http://www.su.edu}{\includegraphics[height=1.75cm]{sulogo.eps}}
\vspace{-1.69cm}

{{\ }\hfill\small
\begin{tabular}{cl}
& Math 207\\
& Introduction to Statistics\\
\end{tabular}
}
\setlength{\baselineskip}{1.05\baselineskip}
\bigskip

\begin{center}
\textbf{\large  Quiz 1}
\end{center}
\medskip

1. A hypothetical population distribution is shown in the table below.

{\ }\hspace{10pt}\begin{tabular}{cc}
       & Percent of\\
Age    & Population\\\hline
0--5   & 10\vphantom{\Large Y}\\  % 2
\Z6--12  & 12\\                   % 2
13--18 & 15\\                     % 3
19--34 & 45\\                     % 4
35--45 & 10\\                     % 1
46--56 & \Z5\\                    % 1/2
57--69 & \Z3                      % 1/4
\end{tabular}
\vspace{.5in}

\hspace{5pt} a) Draw a histogram of the population distribution.  Label
both the axes.

\hspace{5pt} b) About what percent  of the population is between 8 and 11 years old?
\vspace{.5in}

%\hspace{5pt} c) About what percent  of the population is between 44 and 46 years old?
%\vspace{.5in}

2. A fair die is rolled 10 times.  The number of dots showing on the rolls
was:  5, 3, 1, 6, 5, 3, 4, 3, 4, 4. On the back side of this page, 
draw a histogram of this data.
\vspace{.5in}


3. (Hypothetical.)  In a clinical trial, data collection usually starts 
at \lq\lq baseline,\rq\rq\ when the subjects are recruited into the trial 
but before they are assigned to treatment or control.  Data collection continues 
until the end of followup.  Two clinical trials on prevention of heart attacks 
report baseline data on smoking, shown below.  In one of these trials, 
the randomization did not work.  Which one was it and and how can you tell?
\vspace{-12pt}
\begin{center}
\begin{tabular}{cc}
\begin{tabular}{l}
\\
\\
\rput(-12pt,-3pt){(i)}
Treatment\vphantom{\Large Y}\\
Control\\[10pt]
\rput(-12pt,-3pt){(ii)}
Treatment\\
Control
\end{tabular}
&
\begin{tabular}{cc}
Number of & Percent\\
persons & who smoked\\\hline
1,012  & 59.3\%\vphantom{\Large Y}\\
\Zc997  & 59.0\%\\[10pt]
\Zc995  & 49.3\%\\
1,017   & 69.0\%
\end{tabular}
\end{tabular}
\end{center}
\vfill
\eject


\href{http://www.su.edu}{\includegraphics[height=1.75cm]{sulogo.eps}}
\vspace{-1.69cm}

{{\ }\hfill\small
\begin{tabular}{cl}
& Math 207\\
& Introduction to Statistics\\
\end{tabular}
}
\setlength{\baselineskip}{1.05\baselineskip}
\bigskip

\begin{center}
\textbf{\large  Quiz 1}
\end{center}
\medskip

1. (Hypothetical.)  In a clinical trial, data collection usually starts 
at \lq\lq baseline,\rq\rq\ when the subjects are recruited into the trial 
but before they are assigned to treatment or control.  Data collection continues 
until the end of followup.  Two clinical trials on prevention of heart attacks 
report baseline data on smoking, shown below.  In one of these trials, 
the randomization did not work.  Which one was it and and how can you tell?
\vspace{-12pt}
\begin{center}
\begin{tabular}{cc}
\begin{tabular}{l}
\\
\\
\rput(-12pt,-3pt){(i)}
Treatment\vphantom{\Large Y}\\
Control\\[10pt]
\rput(-12pt,-3pt){(ii)}
Treatment\\
Control
\end{tabular}
&
\begin{tabular}{cc}
Number of & Percent\\
persons & who smoked\\\hline
1,012  & 49.3\%\vphantom{\Large Y}\\
\Zc997  & 69.0\%\\[10pt]
\Zc995  & 59.3\%\\
1,017   & 59.0\%
\end{tabular}
\end{tabular}
\end{center}
\vspace{1in}

2. A hypothetical population distribution is shown in the table below.

{\ }\hspace{10pt}\begin{tabular}{cc}
       & Percent of\\
Age    & Population\\\hline
0--5   & 5\vphantom{\Large Y}\\  % width=5, height=1
\Z6--12  & 12\\  % width=6, height=2
13--18 & 10\\  % width=5, height=2
19--34 & 15\\  % width=15, height=1
35--45 & 30\\  % width=10, height=3
46--56 & 20\\ % width=10, height=2
57--73 & \Z8   % width=16, height=1/2
\end{tabular}
\vspace{.5in}

\hspace{5pt} a) Draw a histogram of the population distribution.  Label
both the axes.

\hspace{5pt} b) About what percent of the population is between 41 and 44 years old?
\vspace{.5in}

%\hspace{5pt} c) About what percent  of the population is between 44 and 46 years old?
%\vspace{.5in}

3. A fair die is rolled 10 times.  The number of dots showing on the rolls
was:  5, 6, 4, 5, 3, 3, 1, 5, 2, 2. On the back side of this page, 
draw a histogram of this data.

\vfill
\eject

\end{document}

